\documentclass[a4paper,12pt]{article}

\usepackage{fontspec}
\usepackage{xunicode}
\usepackage{xltxtra}
\usepackage[margin=2.5cm]{geometry}
\usepackage{enumitem}
\usepackage[table]{xcolor}
\usepackage{graphicx}
\usepackage{longtable}
\usepackage{booktabs}
\usepackage{tikz}
\usetikzlibrary{shapes.geometric,arrows,positioning,fit,calc,trees}
\usepackage{float}
\usepackage[normalem]{ulem}
\usepackage{needspace}
\usepackage{amsmath}
\usepackage{amssymb}
\usepackage{multicol}
\usepackage{array}

\setmainfont{TH Sarabun New}
\setmonofont{Courier New}

\XeTeXlinebreaklocale "th"
\XeTeXlinebreakskip = 0pt plus 1pt

\linespread{1.15}

\definecolor{darkblue}{RGB}{0,55,107}
\definecolor{teal}{RGB}{0,123,138}
\definecolor{lightgray}{RGB}{242,242,242}

\newcommand{\blank}[1]{\fbox{\rule{0pt}{1.1em}\hspace{#1}}}
\newcommand{\ans}[1]{\textbf{#1}}
\newcommand{\tf}[1]{\fbox{\rule{0pt}{1.1em}\hspace{2.2cm}}}

\begin{document}

% ─── Cover ──────────────────────────────────────────────────────────────────
\begin{center}
    \textbf{\Large CPE 232 Data Models}\\[0.2cm]
    \textbf{\large Mock Final Exam (ข้อสอบจำลองปลายภาค)}\\[0.15cm]
    เวลาสอบ: 3 ชั่วโมง \quad|\quad คะแนนเต็ม: 100 คะแนน \quad|\quad สอบ Onsite
\end{center}
\hrule
\vspace{0.35cm}

\noindent\textbf{ประกาศและขอบเขตข้อสอบ}\\
ข้อสอบครอบคลุมเนื้อหา Lecture 6--9 ได้แก่ Machine Learning Classification, Regression, Parameter Tuning และ Clustering
ข้อสอบ\textbf{เน้นข้อเขียนและคำนวณ (ประมาณ 90\%)} MCQ มีเพียงเล็กน้อย\\
\textbf{สิ่งที่อนุญาต:} เครื่องเขียน, บัตรนักศึกษา และ \textbf{Formula Sheet} ที่อาจารย์แจกให้ (ไม่ต้องท่องสูตร)\\
\textbf{โครงสร้างคะแนน:} Part 1 (10) \;|\; Part 2 (10) \;|\; Part 3 (15) \;|\; Part 4 (30) \;|\; Part 5 (35)

\vspace{0.4cm}

% ─── Part 1 ─────────────────────────────────────────────────────────────────
\section*{Part 1: Multiple Choice / Multiple Selection / True-False (10 คะแนน)}

\subsection*{1.1 Multiple Choice — เลือกคำตอบที่ถูกต้องที่สุด 1 ข้อ (5 ข้อ ข้อละ 1 คะแนน)}

\begin{enumerate}[label=\textbf{\arabic*.}]

\needspace{7\baselineskip}
\item งานใดต่อไปนี้ที่ไม่เหมาะกับ Clustering?
\begin{enumerate}[label=\alph*.]
    \item การแบ่งกลุ่มลูกค้าตามพฤติกรรมการซื้อ
    \item การสรุปลักษณะประชากรในแต่ละกลุ่ม
    \item การดึงข้อมูลจากฐานข้อมูลด้วยเงื่อนไขที่กำหนด
    \item การหาคนที่มีพฤติกรรมแตกต่างไปจากคนอื่น (Anomaly)
\end{enumerate}

\needspace{7\baselineskip}
\item การวัดประสิทธิผลข้อใดต่อไปนี้ที่ไม่เหมาะสมกับแบบจำลอง Regression?
\begin{enumerate}[label=\alph*.]
    \item RMSE
    \item ^2$ Score
    \item MAE
    \item Precision
\end{enumerate}

\needspace{7\baselineskip}
\item แบบจำลอง Machine Learning ใดที่ใช้หลักการหาความเหมือน (Similarity) เป็นหลักในการทำนาย?
\begin{enumerate}[label=\alph*.]
    \item Logistic Regression
    \item K-Nearest Neighbors (KNN)
    \item Decision Tree
    \item Linear Regression
\end{enumerate}

\needspace{7\baselineskip}
\item สมการ Linear Regression ในข้อใดมี Interaction Term?
\begin{enumerate}[label=\alph*.]
    \item  = c_0 + c_1x_1$
    \item  = c_0 + c_1x_1 + c_2x_2$
    \item  = c_0 + c_1x_1 + c_2x_2 + c_3x_1x_2$
    \item  = c_0 x_1 + c_1 x_2$
\end{enumerate}

\needspace{7\baselineskip}
\item ข้อใดไม่ถูกต้องเกี่ยวกับ Hierarchical Clustering?
\begin{enumerate}[label=\alph*.]
    \item Dendrogram เป็นตัวแทนที่แสดงให้เห็นถึงความใกล้เคียงของข้อมูล
    \item ในขั้นตอนแรก ข้อมูลแต่ละตัวอย่างจะถูกถือว่าเป็นหนึ่ง Cluster
    \item เราสามารถหากลุ่ม (Cluster) ได้โดยการตัด Dendrogram ที่ระดับความสูงที่เหมาะสม
    \item ข้อมูลแต่ละตัวจะถูกกำหนดเข้าไปอยู่ใน Cluster ที่มี Centroid ใกล้ที่สุดในการอัปเดตแต่ละรอบ
\end{enumerate}

\end{enumerate}

\subsection*{1.2 Multiple Selection — เลือกทุกข้อที่ถูกต้อง (3 ข้อ ข้อละ 1 คะแนน)}
\textit{หมายเหตุ: ต้องเลือกครบถ้วนทุกข้อที่ถูกจึงจะได้คะแนน}

\begin{enumerate}[label=\textbf{\arabic*.},start=6]

\needspace{8\baselineskip}
\item ข้อใดบ้างเป็น \textbf{Hyperparameter} (กำหนดก่อน training)?
\begin{enumerate}[label=\alph*.]
    \item \texttt{max\_depth} ของ Decision Tree
    \item ค่า weight ที่ model เรียนรู้ระหว่าง Gradient Descent
    \item Learning Rate ($\alpha$) ใน Gradient Descent
    \item จำนวนคลัสเตอร์ (k) ใน K-Means (\texttt{n\_estimators})
    \item ค่า bias term ที่ optimize ผ่าน loss function
\end{enumerate}

\needspace{8\baselineskip}
\item ข้อใดบ้างช่วย\textbf{ลด Overfitting} ได้?
\begin{enumerate}[label=\alph*.]
    \item เพิ่มข้อมูล training
    \item ลด \texttt{max\_depth} ใน Decision Tree
    \item เพิ่มจำนวน features ใน dataset
    \item ใช้ L1/L2 Regularization
    \item เพิ่ม learning rate เป็น 2 เท่า
\end{enumerate}

\needspace{8\baselineskip}
\item ข้อใดบ้างเป็นข้อความที่\textbf{ถูกต้อง}เกี่ยวกับ Silhouette Score?
\begin{enumerate}[label=\alph*.]
    \item ค่าอยู่ในช่วง $[-1, +1]$
    \item ค่าใกล้ +1 บ่งชี้ว่า cluster มีความแน่นและแยกจากกันชัดเจน
    \item ค่าลบหมายความว่าจุดนั้นถูก assign ผิด cluster
    \item ค่า 0 บ่งชี้ว่า cluster สมบูรณ์แบบ
    \item ใช้ได้เฉพาะกับ K-Means เท่านั้น
\end{enumerate}

\end{enumerate}

\subsection*{1.3 True/False — ระบุ จริง (T) หรือ เท็จ (F) (4 ข้อ ข้อละ 0.5 คะแนน)}
\textit{เขียน T หรือ F ในช่องว่าง}

\begin{enumerate}[label=\textbf{\arabic*.},start=9]
    \item \tf{} \quad Entropy = 0 หมายความว่า node นั้น\textbf{ไม่บริสุทธิ์} (impure) อย่างสมบูรณ์
    \vspace{0.3cm}
    \item \tf{} \quad ค่า $R^2$ (Coefficient of Determination) อาจมีค่าติดลบได้ในบางกรณี
    \vspace{0.3cm}
    \item \tf{} \quad K-Means รับประกันว่า converge ไปที่ global optimum เสมอ
    \vspace{0.3cm}
    \item \tf{} \quad Stratified k-Fold Cross-Validation รักษาสัดส่วน class ในแต่ละ fold ให้ใกล้เคียงกับ dataset ต้นฉบับ
    \vspace{0.3cm}
\end{enumerate}

\newpage

% ─── Part 2 ─────────────────────────────────────────────────────────────────
\section*{Part 2: Definition — เขียนนิยามความหมายของคำศัพท์ต่อไปนี้ (10 คะแนน, ข้อละ 1 คะแนน)}
\textbf{คำชี้แจง:} ตอบเป็นภาษาไทยหรือภาษาอังกฤษก็ได้ ให้นิยามที่ครอบคลุมทั้งความหมายและบริบทการใช้งาน

\begin{enumerate}[label=\textbf{\arabic*.}]
    \item \textbf{Machine Learning} \hfill [Lecture 6]
    \vspace{2.2cm}

    \item \textbf{Supervised Learning} \hfill [Lecture 6]
    \vspace{2.2cm}

    \item \textbf{Unsupervised Learning} \hfill [Lecture 9]
    \vspace{2.2cm}

    \item \textbf{Regularization (L1/L2)} \hfill [Lecture 7/8]
    \vspace{2.2cm}

    \item \textbf{Confusion Matrix} \hfill [Lecture 6]
    \vspace{2.2cm}

    \item \textbf{Overfitting} \hfill [Lecture 8]
    \vspace{2.2cm}

    \item \textbf{Hyperparameter} \hfill [Lecture 8]
    \vspace{2.2cm}

    \item \textbf{k-Fold Cross-Validation} \hfill [Lecture 8]
    \vspace{2.2cm}

    \item \textbf{Silhouette Score} \hfill [Lecture 9]
    \vspace{2.2cm}

    \item \textbf{DBSCAN (Density-Based Spatial Clustering)} \hfill [Lecture 9]
    \vspace{2.2cm}
\end{enumerate}

\newpage

% ─── Part 3 ─────────────────────────────────────────────────────────────────
\section*{Part 3: Short Answer \& Data Interpretation (15 คะแนน, ข้อละ 3 คะแนน)}
\textbf{คำชี้แจง:} ตอบสั้น 2--4 บรรทัด เน้นการ ``ตีความ'' และ ``อธิบายเหตุผล'' ไม่ใช่เพียงแค่นิยาม

\begin{enumerate}[label=\textbf{\arabic*.}]

\needspace{8\baselineskip}
\item \textbf{[Lecture 6]} หากสร้างแบบจำลอง Decision Tree แล้วพบว่ามีปัญหา Overfitting (จำข้อมูล Training ได้เกือบหมด แต่ Test Accuracy ต่ำมาก) เราสามารถปรับ Hyperparameter ใดได้บ้างเพื่อแก้ปัญหานี้? (ยกตัวอย่างมา 2 พารามิเตอร์พร้อมอธิบายพอสังเขป)
\vspace{3.5cm}

\needspace{8\baselineskip}
\item \textbf{[Lecture 8]} ระบบ AI ผ่าน training ด้วย Training Accuracy = 98\% แต่เมื่อทดสอบกับข้อมูลใหม่ Test Accuracy = 62\% ปัญหาคืออะไร? ควรแก้ไขอย่างไรบ้าง (ระบุอย่างน้อย 2 วิธี)
\vspace{3.5cm}

\needspace{10\baselineskip}
\item \textbf{[Lecture 6]} จาก Confusion Matrix ด้านล่าง จงตีความผลลัพธ์และอธิบายว่า classifier นี้มีจุดแข็งและจุดอ่อนอะไร ควรปรับปรุงอย่างไร

\begin{center}
\renewcommand{\arraystretch}{1.3}
\begin{tabular}{c|cc}
    \hline
    \rowcolor{cyan!15}
    & \textbf{ทำนาย: บวก (+)} & \textbf{ทำนาย: ลบ (-)} \\
    \hline
    \textbf{จริง: บวก (+)} & 20 (TP) & 80 (FN) \\
    \textbf{จริง: ลบ (-)} & 5 (FP) & 95 (TN) \\
    \hline
\end{tabular}
\end{center}
\vspace{3.5cm}

\needspace{8\baselineskip}
\item \textbf{[Lecture 9]} ในการทำ Clustering เพราะเหตุใดจึงอาจพิจารณาเลือกใช้ \textbf{K-Medoids} แทน \textbf{K-Means}? อธิบายความแตกต่างของทั้ง 2 วิธีนี้มาพอสังเขป
\vspace{3.5cm}

\needspace{8\baselineskip}
\item \textbf{[Business Application]} สมมติคุณสร้างโมเดลทำนายการซื้อสินค้าและพบผลลัพธ์ว่า `กลุ่มลูกค้าที่น่าจะซื้อสินค้าใหม่มากที่สุดคือ กลุ่มนิสิตนักศึกษา ในช่วงเวลา 20.00-24.00 น.'' ในเชิงธุรกิจ คุณจะออกนโยบายหรือกลยุทธ์ (Take action) อย่างไรเพื่อเพิ่มยอดขายหรือสร้างประสบการณ์ที่ดีให้กับลูกค้ากลุ่มนี้? (ยกตัวอย่าง 2 กลยุทธ์)
\vspace{3.5cm}

\end{enumerate}

\newpage

% ─── Part 4 ─────────────────────────────────────────────────────────────────
\section*{Part 4: Long Answer / Essay (30 คะแนน)}
\textbf{คำชี้แจง:} ตอบอย่างครบถ้วน ใช้ตารางหรือเรียงความก็ได้ ประเมินจากความครบถ้วน ความถูกต้อง และความชัดเจน

\subsection*{ข้อ 1 — การเลือกใช้เทคนิค Machine Learning และ Algorithm (10 คะแนน)}

บริษัท E-Commerce แห่งหนึ่งมีโจทย์ทางธุรกิจ 3 โจทย์ด่านล่าง:
\begin{itemize}
    \item \textbf{โจทย์ A:} ทายว่าลูกค้าจะกดยกเลิกสั่งซื้อสินค้าหรือไม่ (บอกเป็น ใช่ / ไม่ใช่)
    \item \textbf{โจทย์ B:} คาดการณ์ยอดใช้จ่ายของลูกค้าแต่ละคนในเดือนหน้า (มูลค่าเป็นบาท)
    \item \textbf{โจทย์ C:} แบ่งกลุ่มลูกค้าเพื่อนำไปจัดแคมเปญโปรโมชันได้ตรงตามไลฟ์สไตล์
\end{itemize}

\textbf{จงเลือกจับคู่โจทย์ A, B, C กับเทคนิค ML ที่เหมาะสม (Classification, Regression, Clustering)} 
พร้อมทั้งแนะนำ \textbf{Algorithm ที่เรียนมาในวิชานี้} อย่างน้อย 1 Algorithm สำหรับแต่ละโจทย์ และให้เหตุผลสนับสนุนสั้นๆ

\vspace{9.5cm}

\newpage

\subsection*{ข้อ 2 — Types of Cross-Validation (10 คะแนน)}

\textbf{ส่วนที่ 2.1} อธิบาย Cross-Validation 4 แบบต่อไปนี้ โดยระบุ concept, ข้อดี, ข้อจำกัด, และเหมาะกับ dataset แบบใด:
\begin{itemize}
    \item k-Fold Cross-Validation
    \item Stratified k-Fold Cross-Validation
    \item Leave-One-Out Cross-Validation (LOO-CV)
    \item Time Series Cross-Validation
\end{itemize}

\vspace{7cm}

\textbf{ส่วนที่ 2.2} สถานการณ์ต่อไปนี้ควรใช้ Cross-Validation แบบใด พร้อมให้เหตุผล (1 คะแนน/ข้อ)

\begin{enumerate}[label=\alph*.]
    \item Dataset ราคาหุ้น SET50 ย้อนหลัง 5 ปี ต้องการ predict ราคาวันถัดไป\\
    ตอบ: \blank{8cm}
    \vspace{0.5cm}
    \item Dataset ตรวจมะเร็ง 200 samples (Positive 10\%, Negative 90\%)\\
    ตอบ: \blank{8cm}
    \vspace{0.5cm}
\end{enumerate}

\newpage

\subsection*{ข้อ 3 — การประยุกต์ใช้ 4 Types of Analytics (10 คะแนน)}

สมมติว่าคุณเป็นนักวิเคราะห์ข้อมูลของร้านค้าปลีก (Retail Store) บริษัทพบว่ายอดขายสาขาในห้างตกต่ำลง 15\% จงประยุกต์ใช้ 4 Types of Analytics เพื่อนำเสนอโซลูชัน:

\begin{enumerate}
    \item \textbf{Descriptive Analytics} (2.5 คะแนน) - จะดูข้อมูลส่วนใดหรือหา Insight อะไรเพื่อดูสิ่งที่เกิดขึ้น?
    \vspace{2.5cm}
    \item \textbf{Diagnostic Analytics} (2.5 คะแนน) - จะตั้งสมมติฐานหรือวิเคราะห์หาสาเหตุอย่างไร?
    \vspace{2.5cm}
    \item \textbf{Predictive Analytics} (2.5 คะแนน) - จะสร้างโมเดลคาดการณ์เรื่องใดในอนาคตเพื่อช่วยร้านค้า?
    \vspace{2.5cm}
    \item \textbf{Prescriptive Analytics} (2.5 คะแนน) - จะแนะนำ Automation หรือ Policy optimization อะไรให้ระบบทำงานได้เอง?
    \vspace{2.5cm}
\end{enumerate}

\newpage

% ─── Part 5 ─────────────────────────────────────────────────────────────────
\section*{Part 5: Calculation & Structure (35 คะแนน)}
\textbf{คำชี้แจง:} แสดงวิธีทำให้ครบทุกขั้นตอน คำตอบที่ไม่มีวิธีทำจะไม่ได้คะแนน

\vspace{0.3cm}
\subsection*{ข้อ 5.1 — Confusion Matrix Metrics (10 คะแนน)}

ระบบ AI วินิจฉัยโรคเบาหวาน ทดสอบกับ 200 patients ได้ผลดังนี้:

\begin{center}
\renewcommand{\arraystretch}{1.4}
\begin{tabular}{c|cc|c}
    \hline
    \rowcolor{cyan!15}
    & \textbf{ทำนาย: ป่วย (+)} & \textbf{ทำนาย: ไม่ป่วย (-)} & \textbf{รวม} \\
    \hline
    \textbf{จริง: ป่วย (+)} & 60 & 20 & 80 \\
    \textbf{จริง: ไม่ป่วย (-)} & 40 & 80 & 120 \\
    \hline
    \textbf{รวม} & 100 & 100 & 200 \\
    \hline
\end{tabular}
\end{center}

\textbf{5.1.1} (5 คะแนน) คำนวณค่าต่อไปนี้พร้อมแสดงวิธีคิด:
\begin{itemize}
    \item Accuracy
    \item Precision
    \item Recall (Sensitivity)
    \item Specificity
    \item F1-Score
\end{itemize}
\vspace{6cm}

\textbf{5.1.2} (5 คะแนน) จากผลลัพธ์ข้างต้น จงตีความ: ระบบนี้เหมาะสมสำหรับการวินิจฉัยโรคหรือไม่? metric ใดสำคัญที่สุดในบริบทนี้? และควรปรับปรุงอย่างไร?
\vspace{3.5cm}

\newpage

\subsection*{ข้อ 5.2 — Regression: MSE, RMSE, MAE และ R² (10 คะแนน)}

จาก dataset ต่อไปนี้ model Linear Regression ให้ค่า predict ดังนี้:

\begin{center}
\renewcommand{\arraystretch}{1.3}
\begin{tabular}{|c|c|c|}
    \hline
    \rowcolor{cyan!15}
    \textbf{\#} & $y_{\text{actual}}$ & $y_{\text{predicted}}$ \\
    \hline
    1 & 2  & 2  \\
    2 & 4  & 4  \\
    3 & 5  & 6  \\
    4 & 8  & 8  \\
    5 & 11 & 10 \\
    \hline
\end{tabular}
\end{center}

\textbf{5.2.1} (10 คะแนน) คำนวณค่าต่อไปนี้พร้อมแสดงวิธีคิดทุกขั้น:
\begin{itemize}
    \item MSE (Mean Squared Error)
    \item RMSE (Root MSE)
    \item MAE (Mean Absolute Error)
    \item SSE และ $SS_{total}$
    \item $R^2$ (Coefficient of Determination)
\end{itemize}
\vspace{8cm}

\newpage

\subsection*{ข้อ 5.3 — Hierarchical Agglomerative Clustering + Dendrogram + Nested Clusters (15 คะแนน)}
\textbf{ใช้ Complete Linkage:} $d(\{C_i\}, \{C_j\}) = \max_{x \in C_i,\, y \in C_j} d(x,y)$

กำหนด Proximity Matrix ของจุดข้อมูล 5 จุด (A, B, C, D, E):

\begin{center}
\renewcommand{\arraystretch}{1.3}
\begin{tabular}{|c|ccccc|}
    \hline
    \rowcolor{cyan!15}
    & \textbf{A} & \textbf{B} & \textbf{C} & \textbf{D} & \textbf{E} \\
    \hline
    \textbf{A} & 0 & 1 & 5 & 6 & 7 \\
    \textbf{B} & 1 & 0 & 4 & 6 & 8 \\
    \textbf{C} & 5 & 4 & 0 & 2 & 3 \\
    \textbf{D} & 6 & 6 & 2 & 0 & 3 \\
    \textbf{E} & 7 & 8 & 3 & 3 & 0 \\
    \hline
\end{tabular}
\end{center}

\textbf{5.3.1} (6 คะแนน) แสดงขั้นตอน Agglomerative Clustering ทีละ step จนเหลือ 1 cluster:\\
ระบุ: (1) คู่ที่ merge, (2) ระยะทางที่ merge, (3) Proximity Matrix ที่อัปเดตแล้ว ทุก iteration
\vspace{8cm}

\textbf{5.3.2} (5 คะแนน) วาด \textbf{Dendrogram} โดยระบุแกน Y เป็นระยะทาง (distance) และแสดงระดับที่ merge แต่ละครั้งให้ชัดเจน
\vspace{7cm}

\textbf{5.3.3} (2 คะแนน) วาด \textbf{Nested Cluster Diagram} แสดงการ group จุดข้อมูลในแต่ละ level
\vspace{5cm}

\textbf{5.3.4} (2 คะแนน) ถ้าตัด Dendrogram ที่ระดับ $d = 4$ จะได้กี่ cluster และประกอบด้วยจุดใดบ้าง?
\vspace{1.5cm}

\newpage

% ─── ANSWER KEY ─────────────────────────────────────────────────────────────
\section*{เฉลยและแนวคิด (Answer Key)}
\textit{ส่วนนี้สำหรับตรวจหลังทำข้อสอบ แนะนำให้ทำข้อสอบเต็มเวลาก่อนเปิดเฉลย}

\subsection*{Part 1: คำตอบ}

\noindent\textbf{1.1 MCQ}
\begin{multicols}{5}
\begin{enumerate}[label=\textbf{\arabic*.}]
    \item \ans{c.}
    \item \ans{d.}
    \item \ans{b.}
    \item \ans{c.}
    \item \ans{d.}
\end{enumerate}
\end{multicols}

\noindent\textbf{1.2 Multiple Selection}
\begin{itemize}
    \item \textbf{ข้อ 6:} \ans{a, c, d} (max\_depth, Learning Rate, n\_estimators คือ Hyperparameters; b,e คือ learned parameters)
    \item \textbf{ข้อ 7:} \ans{a, b, d} (เพิ่มข้อมูล, ลด max\_depth, L1/L2 Reg ลด overfit; c เพิ่ม features อาจทำ overfit มากขึ้น; e เพิ่ม LR อาจทำให้ diverge)
    \item \textbf{ข้อ 8:} \ans{a, b, c} (ช่วง -1 ถึง +1, ค่าใกล้ +1 ดี, ค่าลบ = misclustered; d ผิด ค่า 0 = on boundary; e ผิด Silhouette ใช้กับ clustering ทั่วไป)
\end{itemize}

\noindent\textbf{1.3 True/False}
\begin{multicols}{4}
\begin{enumerate}[label=\textbf{\arabic*.},start=9]
    \item \ans{F} (Entropy=0 คือ pure)
    \item \ans{T} (R² $<$ 0 ได้)
    \item \ans{F} (แค่ local optima)
    \item \ans{T} (Stratified รักษา class ratio)
\end{enumerate}
\end{multicols}

\subsection*{Part 2: แนวคำตอบ (Definitions)}

\begin{enumerate}[label=\textbf{\arabic*.}]
    \item \textbf{Entropy:} วัดความไม่แน่นอน (impurity) ของ node ใน Decision Tree $H(S) = -\sum_{i} p_i \log_2 p_i$ ค่า 0 = pure (class เดียว), ค่า 1 = equally mixed สองคลาส
    \item \textbf{Information Gain:} การลดลงของ Entropy หลังแบ่งข้อมูลด้วย attribute นั้น $IG(S,A) = H(S) - \sum_v \frac{|S_v|}{|S|}H(S_v)$ ยิ่งสูงยิ่งเป็น attribute ที่ดีสำหรับ split
    \item \textbf{Confusion Matrix:} ตาราง 2×2 (หรือ n×n) แสดงผล classification ประกอบด้วย TP (ทำนายถูก-บวก), TN (ทำนายถูก-ลบ), FP (ทำนายผิด-positive), FN (ทำนายผิด-negative)
    \item \textbf{Overfitting:} สภาวะที่ model เรียนรู้ noise ใน training data จน Training acc สูงมาก แต่ Test acc ต่ำ --- model ไม่สามารถ generalize กับข้อมูลใหม่ได้
    \item \textbf{Gradient Descent:} Algorithm ปรับ parameter โดย update ในทิศทาง\textbf{ตรงข้าม} gradient ของ loss $\theta := \theta - \alpha \cdot \nabla_\theta L(\theta)$ เพื่อลด loss function ทีละขั้น
    \item \textbf{Hyperparameter:} Parameter ที่\textbf{กำหนดก่อน training} และไม่ถูก optimize โดย model เช่น max\_depth, learning\_rate, k ใน KNN ต่างจาก learnable parameters ที่ model เรียนรู้ระหว่าง training
    \item \textbf{k-Fold CV:} แบ่ง dataset เป็น k ส่วนเท่า ๆ กัน วนซ้ำ k รอบ โดยแต่ละรอบใช้ส่วนต่างกันเป็น validation set ที่เหลือเป็น training ค่า metric เฉลี่ยจาก k รอบ = ประสิทธิภาพของ model
    \item \textbf{Centroid:} จุดศูนย์กลางของ cluster ใน K-Means คำนวณจากค่าเฉลี่ยของ feature ทุกมิติของจุดทั้งหมดใน cluster $\mu_k = \frac{1}{|C_k|}\sum_{x \in C_k} x$
    \item \textbf{Silhouette Score:} วัดคุณภาพ clustering ช่วง $[-1, +1]$ โดย $s(i) = \frac{b(i)-a(i)}{\max(a(i),b(i))}$ โดย $a$=ระยะเฉลี่ยใน cluster, $b$=ระยะเฉลี่ยไปยัง cluster ใกล้ที่สุด ค่าสูงหมายถึง cluster ดี
    \item \textbf{DBSCAN:} Algorithm clustering แบบ density-based กำหนด core points ด้วย $\varepsilon$ (รัศมี) และ MinPts (จำนวนจุดขั้นต่ำ) ไม่ต้องระบุ k ล่วงหน้า จัดการ noise ได้ และ cluster มีรูปร่างได้อิสระ
\end{enumerate}

\subsection*{Part 3: แนวคำตอบ}

\begin{enumerate}[label=\textbf{\arabic*.}]
    \item \textbf{Decision Tree Hyperparameters:} เพื่อลด Overfitting สามารถปรับพารามิเตอร์เช่น (1) ลด \texttt{max\_depth} (จำกัดความลึกของต้นไม้ไม่ให้แตกกิ่งลึกเกินไป) หรือ (2) เพิ่ม \texttt{min\_samples\_split/leaf} (กำหนดจำนวน sample ขั้นต่ำในการแตกกิ่ง เพื่อไม่ให้โมเดลสร้างกิ่งสำหรับข้อมูลที่น้อยและเฉพาะเจาะจงเกินไป)

    \item \textbf{Overfitting:} Training 98\%, Test 62\% คือ Overfitting ชัดเจน Model จดจำ training data มากเกินไป วิธีแก้: (1) ลด max\_depth / เพิ่ม min\_samples\_split ใน DT (2) เพิ่ม L1/L2 Regularization (3) ใช้ k-Fold CV เพื่อ estimate performance จริง (4) เพิ่มข้อมูล training (5) Early Stopping

    \item \textbf{CM Interpretation:} TP=20, FN=80 หมายความว่า Recall = 20/100 = 0.20 ต่ำมาก --- model ``พลาด'' ผู้ป่วยจริง 80\% ถือว่า\textbf{อันตรายมาก} ในบริบทการแพทย์ ควรปรับ decision threshold ให้ต่ำลง  เพื่อให้ model ตรวจจับ positive case ได้ดีขึ้น

    \item \textbf{Silhouette 0.12:} ค่า 0.12 ใกล้ 0 บ่งชี้ว่า cluster ทับซ้อนกัน หรือจุดข้อมูลอยู่บน boundary ระหว่าง cluster --- คุณภาพการ cluster\textbf{ต่ำมาก} วิธี diagnose: ใช้ Elbow Method (plot WCSS vs k) หรือ Silhouette Score vs k เพื่อหาค่า k ที่เหมาะสม

    \item \textbf{Local Minimum:} เกิดเมื่อ gradient ลด loss ลงมาจนติดอยู่ที่จุดต่ำสุดเฉพาะที่ที่ไม่ใช่ global minimum gradient = 0 แต่ไม่ใช่คำตอบที่ดีที่สุด Adam optimizer แก้ปัญหาด้วยการรวม Momentum (สะสม gradient) และ Adaptive Learning Rate (ปรับ LR ตาม parameter) ทำให้ออกจาก local minimum และ saddle point ได้ดีกว่า SGD ธรรมดา
\end{enumerate}

\newpage
\subsection*{Part 4: แนวคำตอบ}

\subsubsection*{ข้อ 1: การเลือกใช้เทคนิค Machine Learning}

\textbf{โจทย์ A (Churn/Cancel Prediction):} 
\begin{itemize}
    \item \textbf{เทคนิค:} Classification (เพราะผลลัพธ์เป็นหมวดหมู่ ใช่/ไม่ใช่)
    \item \textbf{Algorithm ยอมรับคำตอบเช่น:} Decision Tree, Logistic Regression, หรือ KNN (ตอบมา 1 อย่าง) โดยให้เหตุผลเช่น Decision Tree สามารถอธิบายเหตุผลด้วยกฎได้ชัดเจน
\end{itemize}

\textbf{โจทย์ B (Spending Prediction):} 
\begin{itemize}
    \item \textbf{เทคนิค:} Regression (เพราะผลลัพธ์เป็นตัวเลขต่อเนื่อง หรือมูลค่าทางการเงิน)
    \item \textbf{Algorithm ยอมรับคำตอบเช่น:} Linear Regression เป็นต้น โดยให้เหตุผลว่าสามารถหาความสัมพันธ์และน้ำหนักระหว่างตัวแปรกับยอดขายได้
\end{itemize}

\textbf{โจทย์ C (Customer Segmentation):} 
\begin{itemize}
    \item \textbf{เทคนิค:} Clustering / Unsupervised Learning (เพราะเป็นการแบ่งกลุ่มโดยไม่มีผลเฉลยล่วงหน้า)
    \item \textbf{Algorithm ยอมรับคำตอบเช่น:} K-Means, K-Medoids, DBSCAN, หรือ Hierarchical Clustering โดยให้เหตุผลเช่น K-Means เหมาะกับการจัดกลุ่มให้มีจำนวนกลุ่มตามที่กำหนดไว้ได้รวดเร็ว
\end{itemize}

\subsubsection*{ข้อ 2: Cross-Validation Types}

\begin{center}
\renewcommand{\arraystretch}{1.25}
\small
\begin{tabular}{|p{2.8cm}|p{3.2cm}|p{3.2cm}|p{3.0cm}|}
    \hline
    \rowcolor{cyan!15}
    \textbf{ประเภท} & \textbf{Concept} & \textbf{ข้อดี / ข้อจำกัด} & \textbf{เหมาะกับ} \\
    \hline
    k-Fold & แบ่ง k ส่วน วน k รอบ & ประเมิน variance ดี / ใช้เวลา k เท่า & Dataset ทั่วไป balanced \\
    \hline
    Stratified k-Fold & k-Fold แต่รักษา class ratio ทุก fold & ดีกับ imbalanced / ใช้เวลาเท่า k-Fold & Class imbalanced \\
    \hline
    LOO-CV & k = n, test ทีละ 1 sample & ใช้ข้อมูลได้สูงสุด / ช้ามาก O(n) & Dataset เล็ก ($n < 100$) \\
    \hline
    Time Series CV & ขยาย training window ทีละ step & ไม่ data leakage จากอนาคต / ต้องเรียง & Time series, sequential data \\
    \hline
\end{tabular}
\end{center}

\textbf{ส่วน 2.2:} (a) Stock price $\rightarrow$ \textbf{Time Series CV} เพราะข้อมูลมี temporal dependency ห้ามสุ่ม\\
(b) Cancer dataset imbalanced $\rightarrow$ \textbf{Stratified k-Fold} เพื่อรักษา class ratio ใน fold

\newpage
\subsection*{Part 5: เฉลยการคำนวณ}

\subsubsection*{5.1 Entropy \& Information Gain}

Dataset: 8 samples, 4 ใช่ (+), 4 ไม่ใช่ (-) $\Rightarrow$ $H(S) = -\frac{4}{8}\log_2\frac{4}{8} - \frac{4}{8}\log_2\frac{4}{8} = -\frac{1}{2}(-1) - \frac{1}{2}(-1) = \mathbf{1.0}$ bit

\textbf{รายได้ (Income):}
\begin{itemize}
    \item สูง (High): rows 1,2,3 $\to$ 3 ใช่, 0 ไม่ใช่ $\Rightarrow H = 0$
    \item ปานกลาง (Medium): rows 4,5 $\to$ 1 ใช่, 1 ไม่ใช่ $\Rightarrow H = 1.0$
    \item ต่ำ (Low): rows 6,7,8 $\to$ 0 ใช่, 3 ไม่ใช่ $\Rightarrow H = 0$
    \item $IG(\text{Income}) = 1.0 - \frac{3}{8}(0) - \frac{2}{8}(1.0) - \frac{3}{8}(0) = 1.0 - 0.25 = \mathbf{0.75}$
\end{itemize}

\textbf{นักศึกษา (Student):}
\begin{itemize}
    \item ใช่: rows 1,2,4,7 $\to$ 3 ใช่, 1 ไม่ใช่ $\Rightarrow H = -\frac{3}{4}\log_2\frac{3}{4} - \frac{1}{4}\log_2\frac{1}{4} = -\frac{3}{4}(1.585-2) - \frac{1}{4}(-2) = \frac{3}{4}(0.415) + 0.5 \approx 0.811$
    \item ไม่ใช่: rows 3,5,6,8 $\to$ 1 ใช่, 3 ไม่ใช่ $\Rightarrow H = 0.811$ (symmetric)
    \item $IG(\text{Student}) = 1.0 - \frac{4}{8}(0.811) - \frac{4}{8}(0.811) = 1.0 - 0.811 = \mathbf{0.189}$
\end{itemize}

\textbf{อายุ (Age):}
\begin{itemize}
    \item น้อยกว่า 30: rows 1,2,5,6 $\to$ 2 ใช่, 2 ไม่ใช่ $\Rightarrow H = 1.0$
    \item 30--50: rows 3,4,7,8 $\to$ 2 ใช่, 2 ไม่ใช่ $\Rightarrow H = 1.0$
    \item $IG(\text{Age}) = 1.0 - \frac{4}{8}(1.0) - \frac{4}{8}(1.0) = \mathbf{0.0}$
\end{itemize}

\textbf{สรุป IG:} Income = 0.75 $>$ Student = 0.189 $>$ Age = 0.0\\
\textbf{Root Node = รายได้ (Income)} เพราะให้ Information Gain สูงสุด

\subsubsection*{5.2 Confusion Matrix}

TP=60, FP=40, FN=20, TN=80, Total=200
\[
\text{Accuracy} = \frac{60+80}{200} = \frac{140}{200} = \mathbf{0.70}
\]
\[
\text{Precision} = \frac{60}{60+40} = \frac{60}{100} = \mathbf{0.60}
\]
\[
\text{Recall} = \frac{60}{60+20} = \frac{60}{80} = \mathbf{0.75}
\]
\[
\text{Specificity} = \frac{80}{80+40} = \frac{80}{120} \approx \mathbf{0.667}
\]
\[
\text{F1} = \frac{2 \times 0.60 \times 0.75}{0.60 + 0.75} = \frac{0.90}{1.35} = \frac{2}{3} \approx \mathbf{0.667}
\]

\textbf{ตีความ:} Recall = 0.75 หมายถึงผู้ป่วยจริง 25\% ถูกพลาด (FN=20) ในบริบทการแพทย์ FN อันตรายกว่า FP เพราะผู้ป่วยไม่ได้รับการรักษา ควรลด threshold เพื่อเพิ่ม Recall แม้ Precision จะลดลงบ้าง

\subsubsection*{5.3 Regression}

errors $(e_i = y_{actual} - y_{pred})$: 0, 0, -1, 0, 1

\[
\text{MSE} = \frac{0^2+0^2+(-1)^2+0^2+1^2}{5} = \frac{2}{5} = \mathbf{0.40}, \quad
\text{RMSE} = \sqrt{0.40} \approx \mathbf{0.632}
\]
\[
\text{MAE} = \frac{0+0+1+0+1}{5} = \mathbf{0.40}
\]
\[
\bar{y} = \frac{2+4+5+8+11}{5} = 6.0
\quad
\text{SSE} = \sum(y_i - \hat{y}_i)^2 = 0+0+1+0+1 = \mathbf{2}
\]
\[
SS_{total} = (2-6)^2+(4-6)^2+(5-6)^2+(8-6)^2+(11-6)^2 = 16+4+1+4+25 = \mathbf{50}
\]
\[
R^2 = 1 - \frac{SSE}{SS_{total}} = 1 - \frac{2}{50} = \mathbf{0.96}
\]

\textbf{Gradient Descent:}
$\theta_0^{new} = 1.0 - 0.1 \times 2.0 = \mathbf{0.80}$,\quad
$\theta_1^{new} = 0.5 - 0.1 \times (-1.5) = 0.5 + 0.15 = \mathbf{0.65}$

\subsubsection*{5.4 K-Means}

\textbf{Iteration 1} ($\mu_1=(1,1)$, $\mu_2=(8,4)$):
\begin{center}
\small
\begin{tabular}{|c|c|c|c|}
    \hline
    \rowcolor{cyan!15}
    จุด & $d(\cdot, \mu_1)$ & $d(\cdot, \mu_2)$ & Cluster \\
    \hline
    A(1,1) & 0 & $\sqrt{58}\approx7.62$ & 1 \\
    B(2,3) & $\sqrt{5}\approx2.24$ & $\sqrt{37}\approx6.08$ & 1 \\
    C(3,2) & $\sqrt{5}\approx2.24$ & $\sqrt{29}\approx5.39$ & 1 \\
    D(6,5) & $\sqrt{41}\approx6.40$ & $\sqrt{5}\approx2.24$ & 2 \\
    E(7,6) & $\sqrt{61}\approx7.81$ & $\sqrt{5}\approx2.24$ & 2 \\
    F(8,4) & $\sqrt{58}\approx7.62$ & 0 & 2 \\
    \hline
\end{tabular}
\end{center}
$\mu_1^{new} = \left(\frac{1+2+3}{3}, \frac{1+3+2}{3}\right) = (2, 2)$, \quad $\mu_2^{new} = \left(\frac{6+7+8}{3}, \frac{5+6+4}{3}\right) = (7, 5)$

\textbf{Iteration 2} ($\mu_1=(2,2)$, $\mu_2=(7,5)$): ทุกจุดยังอยู่ใน cluster เดิม (A,B,C $\to$ Cluster 1; D,E,F $\to$ Cluster 2)\\
Centroids ไม่เปลี่ยน $\Rightarrow$ \textbf{Converged!} ที่ Cluster 1 = \{A,B,C\}, Cluster 2 = \{D,E,F\}

\subsubsection*{5.5 Hierarchical Agglomerative Clustering}

\textbf{Step 1:} min = $d(A,B) = 1$ $\to$ Merge \{A,B\}\\
Complete Linkage: $d(\{AB\},C)=\max(5,4)=5$, $\;d(\{AB\},D)=\max(6,6)=6$, $\;d(\{AB\},E)=\max(7,8)=8$

\begin{center}
\small
\begin{tabular}{|c|cccc|}
    \hline \rowcolor{cyan!15}
    & \{AB\} & C & D & E \\ \hline
    \{AB\} & 0 & 5 & 6 & 8 \\
    C & 5 & 0 & 2 & 3 \\
    D & 6 & 2 & 0 & 3 \\
    E & 8 & 3 & 3 & 0 \\ \hline
\end{tabular}
\end{center}

\textbf{Step 2:} min = $d(C,D) = 2$ $\to$ Merge \{C,D\}\\
$d(\{AB\},\{CD\})=\max(5,6)=6$, $d(\{CD\},E)=\max(3,3)=3$

\begin{center}
\small
\begin{tabular}{|c|ccc|}
    \hline \rowcolor{cyan!15}
    & \{AB\} & \{CD\} & E \\ \hline
    \{AB\} & 0 & 6 & 8 \\
    \{CD\} & 6 & 0 & 3 \\
    E & 8 & 3 & 0 \\ \hline
\end{tabular}
\end{center}

\textbf{Step 3:} min = $d(\{CD\},E) = 3$ $\to$ Merge \{C,D,E\}\\
$d(\{AB\},\{CDE\}) = \max(6,8) = 8$

\textbf{Step 4:} Merge \{AB\} + \{CDE\} ที่ $d = 8$

\vspace{0.3cm}
\textbf{Dendrogram (เฉลย):}

\begin{center}
\begin{tikzpicture}[scale=0.75]
    \draw[->] (0,-0.5) -- (0,9.5) node[above,font=\small] {ระยะทาง};
    \foreach \y in {1,2,3,4,5,6,7,8} {
        \draw (-0.1,\y) -- (0.1,\y);
        \node[left,font=\scriptsize] at (-0.15,\y) {\y};
    }
    \node[below,font=\scriptsize] at (1,0) {A};
    \node[below,font=\scriptsize] at (2,0) {B};
    \node[below,font=\scriptsize] at (4,0) {C};
    \node[below,font=\scriptsize] at (5,0) {D};
    \node[below,font=\scriptsize] at (6.5,0) {E};
    % A+B at d=1
    \draw (1,0) -- (1,1); \draw (2,0) -- (2,1);
    \draw[thick,blue] (1,1) -- (2,1);
    % C+D at d=2
    \draw (4,0) -- (4,2); \draw (5,0) -- (5,2);
    \draw[thick,teal] (4,2) -- (5,2);
    % {CD}+E at d=3
    \draw (4.5,2) -- (4.5,3); \draw (6.5,0) -- (6.5,3);
    \draw[thick,orange] (4.5,3) -- (6.5,3);
    % {AB}+{CDE} at d=8
    \draw (1.5,1) -- (1.5,8); \draw (5.5,3) -- (5.5,8);
    \draw[thick,red] (1.5,8) -- (5.5,8);
    % Labels
    \node[right,font=\scriptsize,blue] at (2,1.1) {d=1};
    \node[right,font=\scriptsize,teal] at (5,2.1) {d=2};
    \node[right,font=\scriptsize,orange] at (6.5,3.1) {d=3};
    \node[right,font=\scriptsize,red] at (5.5,8.1) {d=8};
\end{tikzpicture}
\end{center}

\textbf{Nested Cluster Diagram (เฉลย):}
\begin{center}
\begin{tikzpicture}[scale=0.9]
    \node[circle,fill=darkblue,inner sep=3pt,label=below:A] (A) at (0,0) {};
    \node[circle,fill=darkblue,inner sep=3pt,label=below:B] (B) at (1.5,0) {};
    \node[circle,fill=darkblue,inner sep=3pt,label=below:C] (C) at (4,0) {};
    \node[circle,fill=darkblue,inner sep=3pt,label=below:D] (D) at (5.5,0) {};
    \node[circle,fill=darkblue,inner sep=3pt,label=below:E] (E) at (7,1.2) {};
    % {A,B}
    \draw[dashed,blue,thick,rounded corners] (-0.5,-0.6) rectangle (2,0.7);
    \node[blue,font=\scriptsize] at (0.75,1.0) {\{A,B\}};
    % {C,D}
    \draw[dashed,teal,thick,rounded corners] (3.5,-0.6) rectangle (6,0.7);
    \node[teal,font=\scriptsize] at (4.75,1.0) {\{C,D\}};
    % {C,D,E}
    \draw[dashed,orange,thick,rounded corners] (3.2,-0.9) rectangle (7.6,1.9);
    \node[orange,font=\scriptsize] at (5.4,2.2) {\{C,D,E\}};
    % {A,B,C,D,E}
    \draw[dashed,red,thick,rounded corners] (-0.8,-1.2) rectangle (8.0,2.7);
    \node[red,font=\scriptsize] at (3.5,3.0) {\{A,B,C,D,E\}};
\end{tikzpicture}
\end{center}

\textbf{5.5.4:} ตัดที่ $d=4$ จะได้ \textbf{2 clusters}: \{A,B\} และ \{C,D,E\}\\
(เพราะ \{A,B\} merge กันที่ $d=1<4$ และ \{C,D,E\} merge กันครบที่ $d=3<4$ แต่ merge กัน 2 กลุ่มใหญ่ที่ $d=8>4$)

\newpage

% ─── FORMULA SHEET ───────────────────────────────────────────────────────────
\section*{Formula Sheet (สูตรสำหรับใช้ในห้องสอบ)}
\textit{อาจารย์จะแจก Formula Sheet แผ่นนี้ในห้องสอบ — ไม่ต้องท่องสูตร แต่ต้องเข้าใจวิธีใช้}

\subsection*{Lecture 6: Classification}
\begin{multicols}{2}
\[
H(S) = -\sum_{i=1}^{c} p_i \log_2 p_i
\]
\[
IG(S,A) = H(S) - \sum_{v \in \text{Values}(A)} \frac{|S_v|}{|S|} H(S_v)
\]
\[
\text{Gini}(S) = 1 - \sum_{i=1}^{c} p_i^2
\]
\[
\text{Sigmoid}(z) = \frac{1}{1+e^{-z}}
\]
\columnbreak
\[
\text{Accuracy} = \frac{TP+TN}{TP+TN+FP+FN}
\]
\[
\text{Precision} = \frac{TP}{TP+FP}, \quad \text{Recall} = \frac{TP}{TP+FN}
\]
\[
\text{F1} = \frac{2 \cdot \text{Precision} \cdot \text{Recall}}{\text{Precision}+\text{Recall}}
\]
\[
\text{Specificity} = \frac{TN}{TN+FP}
\]
\end{multicols}

\subsection*{Lecture 7: Regression}
\begin{multicols}{2}
\[
\text{MSE} = \frac{1}{n}\sum_{i=1}^{n}(y_i - \hat{y}_i)^2
\]
\[
\text{RMSE} = \sqrt{\text{MSE}}, \quad \text{MAE} = \frac{1}{n}\sum|y_i - \hat{y}_i|
\]
\[
R^2 = 1 - \frac{SSE}{SS_{total}} = 1 - \frac{\sum(y_i-\hat{y}_i)^2}{\sum(y_i-\bar{y})^2}
\]
\columnbreak
\[
\theta_j := \theta_j - \alpha \cdot \frac{\partial L}{\partial \theta_j}
\]
\[
L_{\text{Ridge}} = L + \lambda\sum\theta_j^2
\]
\[
L_{\text{Lasso}} = L + \lambda\sum|\theta_j|
\]
\end{multicols}

\subsection*{Lecture 9: Clustering}
\begin{multicols}{2}
\[
d_{\text{Euclidean}} = \sqrt{\sum_{k}(x_k - y_k)^2}
\]
\[
d_{\text{Manhattan}} = \sum_{k}|x_k - y_k|
\]
\[
\text{WCSS} = \sum_{k=1}^{K}\sum_{x \in C_k} \|x - \mu_k\|^2
\]
\columnbreak
\[
s(i) = \frac{b(i) - a(i)}{\max(a(i), b(i))}
\]
\[
d_{\text{Complete}}(C_i, C_j) = \max_{x \in C_i, y \in C_j} d(x,y)
\]
\[
d_{\text{Single}}(C_i, C_j) = \min_{x \in C_i, y \in C_j} d(x,y)
\]
\end{multicols}

\newpage

% ─── RECAP ───────────────────────────────────────────────────────────────────
\section*{RECAP: สรุปเนื้อหาจำสอบ (อ่านก่อนสอบ)}

\subsection*{RECAP 1: Classification (Lecture 6)}
\begin{itemize}
    \item \textbf{Decision Tree:} เลือก root ด้วย IG สูงสุด, Entropy=0 คือ pure node, Gini ใช้แทน Entropy ได้
    \item \textbf{KNN:} หา k เพื่อนบ้านใกล้สุด, ช้าตอน predict, ไม่มี training phase, ต้อง scale features
    \item \textbf{Confusion Matrix จำ:} TP=ทำนายถูก-บวก, FP=ทำนายผิด-บวก (False Alarm), FN=พลาด positive (Missed)
    \item ในทางการแพทย์ FN อันตรายกว่า FP $\Rightarrow$ เน้น Recall; ในการกรอง spam เน้น Precision
\end{itemize}

\subsection*{RECAP 2: Regression (Lecture 7)}
\begin{itemize}
    \item \textbf{MSE ลงโทษ outlier มากกว่า MAE} เพราะยกกำลังสอง; RMSE หน่วยเดียวกับ y
    \item \textbf{$R^2 = 1$} = perfect fit; $R^2 = 0$ = ไม่ดีกว่าการใช้ค่าเฉลี่ย; $R^2 < 0$ = แย่กว่า
    \item \textbf{Gradient Descent:} $\theta := \theta - \alpha \cdot \nabla L$ อัปเดต parameter ลดทาง gradient
    \item \textbf{Batch GD:} ใช้ทุก sample/step (stable แต่ช้า) vs \textbf{SGD:} 1 sample/step (เร็วแต่ noisy) vs \textbf{Mini-batch:} กลาง ๆ
    \item \textbf{Ridge (L2):} ลด magnitude weights, \textbf{Lasso (L1):} ทำให้บาง weight = 0 (feature selection)
\end{itemize}

\subsection*{RECAP 3: Parameter Tuning (Lecture 8)}
\begin{itemize}
    \item \textbf{Hyperparameter vs Learnable:} max\_depth, LR = Hyperparameter (กำหนดก่อน); weights, bias = Learnable (model เรียนรู้)
    \item \textbf{Overfitting:} Train high, Test low $\Rightarrow$ แก้ด้วย: regularization, ลด complexity, เพิ่มข้อมูล, k-Fold CV
    \item \textbf{k-Fold:} ทั่วไป; \textbf{Stratified k-Fold:} imbalanced class; \textbf{LOO-CV:} dataset เล็ก; \textbf{Time Series CV:} sequential
    \item \textbf{Grid Search:} ลอง combination ทั้งหมด (exhaustive แต่ช้า); \textbf{Random Search:} สุ่ม (เร็วกว่า); \textbf{Bayesian:} ใช้ prior ชี้นำ
    \item \textbf{Adam:} รวม Momentum + Adaptive LR ใช้ได้ดีกับส่วนใหญ่ ออกจาก local minimum ได้ดี
    \item \textbf{Learning Rate Scheduling:} ลด LR ตามเวลา (decay) ช่วย converge ได้ละเอียดขึ้น
\end{itemize}

\subsection*{RECAP 4: Clustering (Lecture 9)}
\begin{itemize}
    \item \textbf{K-Means:} (1) init centroids, (2) assign points, (3) update centroids, (4) repeat จน converge WCSS วัดความแน่น
    \item \textbf{Elbow Method:} plot WCSS vs k หา ``ข้อศอก'' ที่ WCSS ลดช้าลงกะทันหัน
    \item \textbf{Silhouette:} $[-1,+1]$; $s \approx 1$ = ดี; $s \approx 0$ = on boundary; $s < 0$ = misclustered
    \item \textbf{Hierarchical Agglomerative:} Bottom-up, Single/Complete/Average Linkage, ผล = Dendrogram ตัด Dendrogram ที่ระดับต่าง ๆ ได้จำนวน cluster ต่าง ๆ
    \item \textbf{Complete Linkage = MAX distance} ระหว่างคู่จุดของ 2 clusters; Single = MIN
    \item \textbf{DBSCAN:} ไม่ต้องกำหนด k, จัดการ noise ได้, cluster ไม่เป็น convex ได้ $\varepsilon$ = รัศมี, MinPts = จุดขั้นต่ำ
    \item \textbf{K-Means ดีเมื่อ:} cluster กลม ๆ, รู้จำนวน k; \textbf{DBSCAN ดีเมื่อ:} cluster รูปร่างซับซ้อน, มี noise/outlier
\end{itemize}

\end{document}
